\documentclass{article}
\usepackage[utf8]{inputenc}
\usepackage[T1]{fontenc}
\usepackage{geometry}
\geometry{
    a4paper,
    left=2.5cm,
    right=2.5cm,
    top=2.5cm,
    bottom=2.5cm
}
\usepackage{graphicx}
\usepackage{float}
\usepackage{hyperref}
\usepackage[english]{babel}
\usepackage{amsmath}
\usepackage{amssymb}
\usepackage{array}
\usepackage{booktabs}
\usepackage{caption}
\usepackage{subcaption}
\usepackage{enumitem}
\usepackage{fancyhdr}
\usepackage{titlesec}
\usepackage{url}
\usepackage{textcomp}
\usepackage{xcolor}
\usepackage{listings}
\usepackage{multirow}
\usepackage{longtable}
\usepackage{tabularx}
\usepackage{times}

% Set up page style
\pagestyle{fancy}
\fancyhf{}
\rfoot{\thepage}
\renewcommand{\headrulewidth}{0pt}

% Title formatting
\titleformat{\chapter}[display]{\normalfont\bfseries}{\chaptertitle\thechapter}{0pt}{\Huge}
\titlespacing*{\chapter}{0pt}{50pt}{40pt}

% Document info
\title{Template Document}
\author{Academic Support Agent}
\date{\today}

\begin{document}
\maketitle

\begin{titlepage}
\begin{center}
\textbf{\LARGE Funda o Getulio Vargas}\\ 
\textbf{\LARGE Escola de Matem tica Aplicada}\\
\textbf{\LARGE Curso de Gradua o em Matem tica Aplicada}

\par
\vspace{170pt}
\textbf{\Large T tulo da disserta o}\\
\vspace{80pt}
\textbf{\Large por Nome}\\
\end{center}

\par
\vfill
\begin{center}
{{\normalsize Rio de Janeiro - Brasil}\\
{\normalsize \the\year}}
\end{center}
\end{titlepage}

\thispagestyle{empty}

\newpage
\begin{center}
\textbf{\LARGE Funda o Getulio Vargas}\\ 
\textbf{\LARGE Escola de Matem tica Aplicada}\\
\textbf{\LARGE Curso de Gradua o em Matem tica Aplicada}

\par
\vspace{100pt}
\textbf{\Large T tulo da disserta o}


\par
\vspace{65pt}
``\textbf{Declaro ser o  nico autor do presente projeto de
monografia que refere-se ao plano de trabalho a ser executado para continuidade da monografia e ressalto que n o recorri a qualquer forma de colabora o ou aux lio de terceiros para realiz -lo a n o ser nos casos e para os fins autorizados pelo professor orientador.}''
\end{center}

\par
\vspace{65pt}
\begin{center}


\hrulefill

\vspace{5pt}
\textbf{\Large Nome}
\end{center}

\par
\vfill
\begin{center}
{{\normalsize Rio de Janeiro - Brasil}\\
{\normalsize \the\year}}
\end{center}

\thispagestyle{empty}

\newpage
\begin{center}
\textbf{\LARGE Funda o Getulio Vargas}\\ 
\textbf{\LARGE Escola de Matem tica Aplicada}\\
\textbf{\LARGE Curso de Gradua o em Matem tica Aplicada}

\par
\vspace{100pt}
\textbf{\Large T tulo da disserta o}


\par
\vspace{65pt}

``\textbf{Projeto de Monografia apresentado   Escola de Matem tica Aplicada como requisito parcial para continuidade ao trabalho de monografia.}''
\end{center}

\par
\vspace{65pt}
\begin{center}

\textbf{Aprovado em } \makebox[30pt]{\hrulefill}\textbf{ de }\makebox[120pt]{\hrulefill}\textbf{ de }\makebox[50pt]{\hrulefill}
\\
\vspace{5pt}
\textbf{Grau atribu do ao Projeto de Monografia:} \makebox[30pt]{\hrulefill}\\
\end{center}


\par
\vspace{40pt}
\begin{center}

\hrulefill

\vspace{5pt}
\textbf{Professor Orientador: }\\
\textbf{Escola de Matem tica Aplicada}\\
\textbf{Funda o Get lio Vargas}
\end{center}

\thispagestyle{empty}


\newpage

\thispagestyle{empty}

\newpage
\section{Introdu o}

introdu o...

\newpage
\section{Objetivo Final}

\subsection{xx}

referencial te rico... \footnote{Ver xx}.

\subsubsection{yy}

\newpage
\subsection{xx}

xx

\newpage
\section{Metodologia}

metodologia...

\subsection{xx}

xx

\newpage
\begin{figure}[H]
\centering
%left bottom right top
%\includegraphics[trim=10mm 00mm 00mm 00mm, scale=0.40]{tabela.pdf}
\begin{center}
Figura 1: tabela.
\end{center}
\end{figure}

\begin{figure}[H]
\centering
%left bottom right top
%\includegraphics[trim=10mm 00mm 00mm 00mm, scale=0.40]{tabela.pdf}
\begin{center}
Figura 2: tabela.
\end{center}
\end{figure}

\newpage
\subsubsection{yy}

yy

% \[\left\{ 
\begin{align*}
f_{i}(x) = (10x + 100), \qquad \text{(1)} \tag{1}\\
f_{ii}(x) = (20x + 200), \qquad \text{(2)} \tag{2} \\
f_{iii}(x) = (30x + 300), \qquad \text{(3)} \tag{3} \\
\end{align*} %\right.\]

xx

\begin{align*}
Vm_{i}(p,l) = ((-1.9141)p + 49.466)l + ((199.51)p - 10795.0), \text {$l$=0} \tag{4}
\end{align*}

\begin{align*}
f_{n}(y) = \frac{y}{1000}, \tag{5}
\end{align*}

\subsection{xx}

xx

%\[\left\{
\begin{align*} 
Funcao_{i}(p) = \gamma + \delta p + \theta p^2 + \omega p^3, \qquad \qquad \qquad \qquad \qquad \qquad \qquad \qquad \text{(6)} \tag{6} \\
\end{align*} %\right.\]

\newpage
\section{Resultados Esperados}

Nesta se o ser o apresentados os resultados esperados...

\subsection{xx}

xx

\newpage
\section{Refer ncias}

[1] a.

\noindent [2] b.

\noindent [3] c.

\noindent [4] d.
\end{document}